\section{结论}

在本报告中，我们介绍了 GLM-5：一款新一代基础模型，实质性弥合了高性能推理与极致计算效率之间的鸿沟。通过从 ``vibe coding'' 范式迈向真正的 ``agentic engineering''，GLM-5 证明了开源权重模型如今已能在复杂真实工作流中与顶级闭源系统同台竞争。
GLM-5 代表了实用 AI 价值的范式转变。通过开源该模型，我们希望赋能社区超越静态基准，探索高效、智能体化通用智能前沿，推动一个由 AI 智能体自主规划、实现并迭代复杂任务的新时代。


\section{彩蛋}
\textbf{``Pony Alpha''} 实验对我们确实是一个关键时刻。将 GLM-5 匿名发布到 OpenRouter 是一次大胆决策，但结果极具说服力。通过去掉品牌名，我们让模型能力本身发声，确保收到的反馈更纯粹、偏差更小。
简要总结如下：


在短短几天内，Pony Alpha 就成为热议焦点。OpenRouter 社区开发者开始注意到其卓越表现，尤其是在复杂编码任务、智能体工作流与角色扮演场景中。

社区猜测此起彼伏，许多用户认为它可能是 Anthropic（Claude Sonnet 5）的泄露更新、神秘 Grok 发布，或 DeepSeek V4。初步统计显示：25\% 用户猜测是 Claude Sonnet 5，20\% 猜测 DeepSeek，10\% 猜测 Grok，其余猜测为 GLM-5。

最终确认它确实是我们的 GLM-5，这一刻意义重大，也有效回应了“中文大模型能否参与前沿竞争”的质疑。
%
Pony Alpha（GLM-5）的成功不仅是基准分数的提升，更意味着我们将关注点转向了工程级可靠性。

这次匿名发布让我们跨越了地缘偏见。社区接纳这个模型，是因为它真的可用。
%
尽管我们庆祝这一成果，也必须保持务实。开源权重模型与最前沿闭源系统的差距正在缩小，但竞赛远未结束。我们仍将坚定推进可扩展、高效率、强智能系统的边界。
